% AdeSEG methodology supplement, 2026-11-20.
% Suggested placement: after baseline benchmark description in content/methodology.
% Evidence labels MUST remain: measured reference vs proposed enhancements.
% This file is a LaTeX fragment, not a standalone document.

\section{Đề xuất mở rộng Lean-SAM2 cho phân đoạn video nội soi}
\label{sec:adeseg-proposed-framework}

\subsection{Phạm vi nghiên cứu và phân biệt bằng chứng}
\label{subsec:adeseg-evidence}

Mục tiêu dài hạn của AdeSEG là phân đoạn riêng biệt mô adenoid và vùng đường
thở vòm họng trong chuỗi video nội soi. Ở giai đoạn hiện tại, PolypGen
\cite{ali2021polypgen} được sử dụng làm bộ dữ liệu \textit{proxy} từ nội soi
đại tràng để kiểm tra các vấn đề về khởi tạo prompt, phân đoạn và lan truyền
theo thời gian. Các kết quả trên PolypGen chưa phải là bằng chứng lâm sàng
cho phân đoạn adenoid hoặc đánh giá mức độ phì đại VA.

Một kết quả tham chiếu đã được lưu trữ là biến thể Lean-SAM2 dựa trên
SAM2.1 Hiera-B+, với bản hiệu chỉnh tương thích indexed spatial RoPE trên
MPS. Phương pháp lịch sử này sử dụng một mask khởi tạo tự động được lưu
trữ cho mỗi video, không tiếp tục phát hiện lại ở các frame sau.
Trên \textbf{21/23 sequence đủ điều kiện khởi tạo}, Dice trung bình đều
theo sequence trên các frame \emph{có foreground} sau khởi tạo là
\textbf{0.742071}, trong khi Dice theo trọng số frame là \textbf{0.772495}
trên 1\,683 frame có foreground. Phương pháp tạo ra \textbf{188 dự đoán
dương tính sai trên 336 frame ground-truth rỗng}. Seq1 và seq7 không có
khởi tạo tự động nên không thuộc mẫu số đánh giá phân đoạn. Seq23 có sử
dụng crop được lựa chọn hồi cứu; các sequence đã từng được xem xét
trước đó. Do vậy, đây là kết quả thăm dò, không phải kiểm định độc lập
trên bộ test chưa từng quan sát.

Nguồn đối chiếu trực tiếp:
\url{https://github.com/aCoderChild/AdeSEG/blob/main/experiments/reference_baselines/modified_lean_742.md}.
Quan trọng: \textbf{0.742071 không phải kết quả của FGPI, RMFA hay TTMR}.
Ba thành phần dưới đây là kiến trúc mở rộng được đề xuất hoặc được cài đặt
thử nghiệm; chưa có thực nghiệm ghép nối trên cùng một backbone Lean-SAM2
để quy kết mức cải thiện cho từng thành phần.

\subsection{Mô hình hóa bài toán và pipeline}
\label{subsec:adeseg-pipeline}

Cho chuỗi ảnh đầu vào $\mathbf{X}=\{X_t\}_{t=0}^{T}$, một vùng mục tiêu
cần phân đoạn được biểu diễn bằng mask nhị phân $Y_t\in\{0,1\}^{H\times W}$.
Trong quá trình inference, hệ thống không truy cập $Y_t$; mask ground-truth
chỉ được dùng cho huấn luyện (đúng split) hoặc đánh giá sau khi dự đoán
đã được cố định. Đầu ra gồm logits $Z_t$ và
$\widehat{Y}_t=\mathbb{I}[\sigma(Z_t)>0.5]$.

Khung phương pháp nghiên cứu gồm ba bài toán liên tiếp:
\textbf{(i)} khởi tạo target bằng detector và đặc trưng ảnh;
\textbf{(ii)} tinh chỉnh biên và phân biệt foreground/background trên
feature map đã được điều kiện hóa theo memory; và
\textbf{(iii)} duy trì thông tin target trong bộ nhớ theo thời gian.
Backbone tham chiếu dự kiến là Lean-SAM2 \footnote{\url{https://arxiv.org/abs/2607.19811}},
vốn kế thừa mô hình promptable video segmentation SAM2
\cite{ravi2024sam2}. Các bước mở rộng được phân tách bằng các
\textit{interface} về prompt, mask logits và memory token.

\paragraph{Giới hạn cài đặt hiện tại.}
FGPI đang được tích hợp với nhánh thử nghiệm MedSAM2 Hiera-T/512; RMFA
và bộ lựa chọn TTMR thử nghiệm dựa trên Lean-SAM2/SAM2.1 Hiera-B+.
Chưa thể mô tả chúng như một mô hình ba module đã chạy đầy đủ trên
cùng backbone hoặc đã tái lập Dice 0.742071.

\subsection{FGPI: Khởi tạo prompt kết hợp detector và anchor đặc trưng}
\label{subsec:adeseg-fgpi}

Tại frame khởi tạo $t_0$, bộ phát hiện YOLO dự đoán bounding box
$b_0=(x_{\min},y_{\min},x_{\max},y_{\max})$. Đặc trưng không gian sớm
$E_0\in\mathbb{R}^{C\times H'\times W'}$ được trích xuất từ image
encoder đã cố định. Tập token trong vùng lõi $C(b_0)$ được xác định
bằng cách thu nhỏ box, còn tập nền $N(b_0)$ được xác định bên ngoài
biên box có khoảng an toàn. Đây là các nhãn giả định theo hình học,
\emph{không phải foreground/background ground truth}.

Sau khi chuẩn hóa token $e_p$, xây dựng bộ
\emph{pseudo-target anchor} $\mathcal{A}_0=\{a_1,\ldots,a_K\}$
bằng cách ưu tiên token có hỗ trợ ở trung tâm và chọn tiếp các
anchor khác nhau về đặc trưng trong vùng ứng viên có hỗ trợ.
Một phiên bản khái quát của tiêu chí chọn anchor tiếp theo là
\begin{equation}
a_{k+1}
=
\arg\max_{p\in C_s\setminus\mathcal{A}_k}
\left[
\alpha\left(1-\max_{a\in\mathcal{A}_k}
\cos(e_p,a)\right)+(1-\alpha)c(p)
\right],
\label{eq:adeseg-anchor-diversity}
\end{equation}
trong đó $c(p)$ biểu diễn prior khoảng cách tới tâm box và $C_s$
là tập ứng viên lõi đã được sàng lọc bằng hỗ trợ không gian.
Công thức này kế thừa \emph{ý tưởng} anchor đa dạng từ Lean-SAM2,
nhưng anchor khởi tạo ở đây chưa được xác nhận bằng mask thật.

Điểm tương phản target và nền được mô hình hóa bằng
\begin{equation}
s(p)=\max_{a\in\mathcal{A}_0}\cos(e_p,a)
-\cos(e_p,\mu_{N}),
\qquad
\mu_{N}=\operatorname{Normalize}
\left(\frac{1}{|N|}\sum_{q\in N}e_q\right).
\label{eq:adeseg-fgpi-score}
\end{equation}
Điểm này cùng với làm trơn không gian và ràng buộc khoảng cách
được dùng để chọn tập positive point $P_0^+$ trong lõi box và
negative point $P_0^-$ bên ngoài biên an toàn. Prompt tổng hợp là
\begin{equation}
\mathcal{P}_0=\{b_0,P_0^+,P_0^-\}.
\label{eq:adeseg-joint-prompt}
\end{equation}
FGPI cung cấp \emph{đồng thời} box và các điểm đặc trưng cho
prompt encoder. Trong cấu hình đề xuất không áp dụng ngưỡng độ
tin cậy của feature để loại bỏ các điểm; nếu hình học đầu vào
không hợp lệ, hệ thống ghi nhận thất bại thay vì ngầm chuyển sang
chỉ dùng YOLO box. Baseline đối chứng vẫn là \emph{YOLO box only}.

\paragraph{Giả thuyết có thể kiểm chứng.}
So với cùng box detector và cùng ngân sách prompt, lựa chọn điểm
dựa trên pseudo-target anchor có thể cải thiện Dice/IoU ở
\emph{frame khởi tạo}. Điều này cần so với cả box-only,
box cộng điểm hình học và box cộng điểm ngẫu nhiên; chưa có
kết quả thực nghiệm chứng minh giả thuyết.

\subsection{RMFA: Tích hợp đặc trưng vùng và tinh chỉnh residual}
\label{subsec:adeseg-rmfa}

RMFA nhận đặc trưng sớm $E_t$, đặc trưng sâu đã được điều kiện hóa
bởi memory $H_t$ và logits gốc $Z_t^{(0)}$ từ \emph{cùng}
backbone/mask decoder. Sau bước chiếu kênh và căn chỉnh kích thước:
\begin{equation}
U_t = \operatorname{Concat}
\big(\operatorname{Up}(W_e E_t),
     \operatorname{Up}(W_h H_t)\big).
\label{eq:adeseg-rmfa-project}
\end{equation}
Xác suất foreground mềm được lấy từ \emph{dự đoán của mô hình},
$p_t=\operatorname{Pool}(\sigma(Z_t^{(0)}))$.
Với $Q_t=W_Q\operatorname{Pool}(U_t)$,
$K_t=W_K\operatorname{Pool}(U_t)$ và
$V_t=W_V\operatorname{Pool}(U_t)$, quan hệ giữa các kênh được
tính riêng cho hai vùng:
\begin{align}
R_{\mathrm{fg}} &=
  \operatorname{Norm}(Q_t\odot p_t)
  \operatorname{Norm}(K_t\odot p_t)^{\top},\\
R_{\mathrm{bg}} &=
  \operatorname{Norm}(Q_t\odot (1-p_t))
  \operatorname{Norm}(K_t\odot (1-p_t))^{\top},\\
A_t &= \operatorname{Softmax}_{\mathrm{channel}}
  (R_{\mathrm{fg}}-R_{\mathrm{bg}}).
\label{eq:adeseg-rmfa-channel}
\end{align}
Tích $A_tV_t$ biểu diễn phép trộn thông tin giữa các kênh
với điều kiện foreground/background. Một nhánh contrast không gian
được tạo từ prototype mềm của hai vùng, rồi kết hợp với feature
fusion và logits gốc để dự đoán correction $\Delta Z_t$:
\begin{equation}
Z_t=Z_t^{(0)}+\Delta Z_t,\qquad
\Delta Z_t =
g_{\phi}\big(
 \operatorname{Concat}(
 U_t+\operatorname{Up}(W_o A_tV_t),
 C_t,Z_t^{(0)})
\big).
\label{eq:adeseg-rmfa-residual}
\end{equation}
Trọng số decoder gốc vẫn được giữ nguyên; lớp cuối của nhánh
residual được khởi tạo bằng không nên tại thời điểm khởi tạo
$Z_t=Z_t^{(0)}$. Đây là tính chất \emph{identity khi khởi tạo},
không phải đảm bảo sai số luôn giảm sau huấn luyện.

Ở nhánh thử nghiệm, loss huấn luyện gồm BCE có trọng số theo biên,
IoU có trọng số và soft Dice; có thể thêm penalty bảo toàn
dự đoán dương tính từ mô hình gốc. Loss và mask thật chỉ dùng
trên TRAIN capture được kiểm soát nguồn dữ liệu. Để giảm nguy cơ
đưa sai số vào memory, phương án nghiên cứu ưu tiên
\textit{output-only refinement}; dẫn logits đã tinh chỉnh vào
memory encoder chỉ được xem là một ablation riêng.

\subsection{TTMR: Nghiên cứu lựa chọn memory theo target và thời gian}
\label{subsec:adeseg-ttmr}

Cần phân biệt TTMR được đề xuất với cơ chế sẵn có của Lean-SAM2:
TAMP đã chọn token bằng target-anchor, TCIM đã gộp memory lịch sử
và duy trì insurance memory, còn TARR đã điều khiển sử dụng
image encoder. TTMR không thể được tuyên bố là đã phát minh
các nguyên lý đó.

Một mở rộng đang nghiên cứu là xếp hạng các memory token lịch sử
$m_i$ tại frame $t_i<t$ với điều kiện target anchor khởi tạo
$\mathcal{A}_0$, điểm sử dụng bộ nhớ gốc $u_i$ và khoảng cách thời gian:
\begin{equation}
q_i =
\lambda_a\max_{a\in\mathcal{A}_0}\cos(m_i,a)
+\lambda_u u_i
+\lambda_r \exp\left(-\frac{t-t_i}{\tau}\right).
\label{eq:adeseg-ttmr-priority}
\end{equation}
Để tránh chọn nhiều token gần như trùng nhau, lựa chọn tham lam
tiếp theo theo
\begin{equation}
i^*=\arg\max_{i\notin S}
\left[q_i-\lambda_d\max_{j\in S}
\max(0,\cos(m_i,m_j))\right],
\quad |S|=B,
\quad t_i<t.
\label{eq:adeseg-ttmr-diversity}
\end{equation}
Các ràng buộc kèm theo gồm ngân sách token cố định $B$,
chỉ dùng token quá khứ, chỉ số duy nhất và độ phủ thời gian
theo frame khi khả thi. Việc so sánh phải giữ \emph{cùng ngân
sách memory} với Lean-SAM2 gốc. Công thức không chứng minh
rằng temporal drift giảm hoặc FPS tăng.

Hiện có mã tham chiếu \textit{offline/shadow} cho quy tắc TTMR,
nhưng chưa có thực nghiệm end-to-end xác nhận việc thay thế
cơ chế memory sống của Lean-SAM2. Do đó TTMR là thành phần
\emph{đề xuất cần đánh giá}, không phải nguồn tạo ra Dice 0.742071.

\subsection{Thiết kế kiểm chứng và phạm vi kết luận}
\label{subsec:adeseg-validity}

Do các mô hình lịch sử và đường cài đặt hiện tại chưa dùng chung
backbone, prompt và protocol, không thể sử dụng chênh lệch Dice
giữa chúng để ước lượng hiệu quả nhân quả của ba module. Khi có
điều kiện chạy thêm, thiết kế ablation tối thiểu gồm:
(i) cùng Lean-SAM2 với box-only;
(ii) box cộng các điểm được chọn hình học;
(iii) box cộng FGPI điểm từ feature;
(iv) (iii) cộng RMFA output-only đã huấn luyện;
(v) kiểm tra memory routing; và
(vi) TTMR với ngân sách memory đúng bằng baseline.

Chỉ số chính của FGPI là Dice/IoU \emph{ngay trên frame khởi tạo}.
Chỉ số của RMFA gồm Dice, IoU và boundary F trên frame được gán nhãn;
TTMR cần thêm Dice theo thời gian, false alarm trên frame rỗng và
chi phí tính toán/bộ nhớ. Cần báo cáo riêng
\emph{frame-weighted mean} và \emph{equal-sequence mean},
cùng số sequence không có prompt và tiêu chuẩn loại trừ.
Độ tin cậy thống kê phải tính theo nhóm nguồn/patient
được xác thực, không bootstrap độc lập từng frame có tương quan.
Temporal IoU cao không đồng nghĩa với phân đoạn chính xác
khi camera di chuyển hoặc mask bị drift ổn định.

\paragraph{Giới hạn cần nêu trong báo cáo.}
Kết quả Lean-SAM2 0.742071 mang tính thăm dò và có crop hồi
cứu; FGPI được thử nghiệm ở nhánh MedSAM2 khác backbone;
RMFA chưa có đánh giá video đầy đủ với adapter học được;
TTMR mới ở mức selector tham chiếu. Các kết quả đó chưa xác lập
tính mới được kiểm chứng, sự vượt trội của kiến trúc ba module
hoặc khả năng ứng dụng lâm sàng trên nội soi vòm họng.
